\section*{Capítulo \ref{ch:b3:planetary-atmospheres} --- Atmosferas planetárias e habitabilidade}

\begin{solution}{exo:b3:planetary-atmospheres:1}
$H = RT/Mg$: Vênus $8.314\times737/(0.04401\times8.87) = \qty{15.7}{km}$;
a Terra $8.314\times288/(0.02897\times9.80) = \qty{8.4}{km}$; Marte
$8.314\times214/(0.04334\times3.71) = \qty{11.1}{km}$; Titã
$8.314\times94/(0.0280\times1.35) = \qty{20.7}{km}$. A atmosfera de Titã,
apesar de fria, é a mais extensa: sua gravidade é sete vezes mais fraca
que a da Terra.
\end{solution}

\begin{solution}{exo:b3:planetary-atmospheres:2}
$\lambda = GMm/kTr = \num{6.674e-11}\times\num{5.97e24}\times\num{1.674e-27}
/(\num{1.381e-23}\,T\times\num{6.878e6})$: 8.78 a \qty{800}{K}, 4.39 a
\qty{1600}{K}. As frações de escape são $9.78\,e^{-8.78} =
\num{1.5e-3}$ e $5.39\,e^{-4.39} = \num{6.7e-2}$: o ciclo solar
altera a fração de escape do hidrogênio por um fator de cerca de 44.
\end{solution}

\begin{solution}{exo:b3:planetary-atmospheres:3}
(a)~$d = S_\mathrm{eff}^{-1/2}$: \omterm{def:b3:planetary-atmospheres:runaway}{estufa úmida} \qty{0.993}{au}, \omterm{def:b3:planetary-atmospheres:runaway}{estufa descontrolada}
\qty{0.975}{au}, estufa máxima \qty{1.705}{au}. (b)~Vênus, a 1.91,
recebe mais que o fluxo da \omterm{def:b3:planetary-atmospheres:runaway}{estufa descontrolada} e ainda mais que o limite empírico
do Vênus recente (1.78): ele está dentro da borda interna. Marte, a 0.43,
situa-se entre 1.014 e 0.3438, dentro da zona habitável conservadora:
sua aridez vem de seu tamanho pequeno e da perda de sua atmosfera, e não
da sua distância.
\end{solution}

\begin{solution}{exo:b3:planetary-atmospheres:4}
(a)~$\log_{10}p = 3.9895 - 443.028/93.66 = -0.741$, $p =
\qty{0.18}{bar}$. (b)~$0.05\times1.6 = \qty{0.080}{bar}$, uma umidade relativa
de $0.080/0.18 = 44\,\%$. (c)~O ar acima dos mares está longe de
saturado: o metano evapora deles, como a água evapora dos
oceanos da Terra, e o ar úmido precisa subir e esfriar antes de chover.
\end{solution}

\begin{solution}{exo:b3:planetary-atmospheres:5}
$\lambda_\mathrm D = 7.02\times2.014/1.008 = 14.0$. $u_\mathrm D/u_\mathrm H
= \sqrt{1.008/2.014} = 0.71$. Razão por átomo $= 0.71\times(15.0\,
e^{-14.0})/(8.02\,e^{-7.02}) = 0.71\times\num{1.69e-3} = \num{1.2e-3}$.
O \omterm{def:b3:planetary-atmospheres:jeans-escape}{escape de Jeans} deixa escapar quase nenhum deutério: é um filtro forte, com
$\alpha \approx 0.001$.
\end{solution}

\begin{solution}{exo:b3:planetary-atmospheres:6}
(a)~$0.016/\num{1.5575e-4} = 103$. (b)~Com $\alpha = 0$, o enriquecimento é
$M_0/M$: $M_0 = 103\times\num{6e15} = \qty{6.2e17}{kg}$, ou
$\num{6.2e17}/(1000\times\num{4.60e14}) = \qty{1.3}{m}$ de água sobre
Vênus. (c)~$M_0/M = 103^{1/0.2} = \num{1.2e10}$, um inventário passado de
\qty{7e25}{kg}, doze vezes a massa da Terra: impossível. O
enriquecimento requer um escape fortemente fracionador (de Jeans ou não térmico),
ativo pelo menos durante a perda da última parte da água, ou um
estado estacionário com suprimento de água.
\end{solution}

\begin{solution}{exo:b3:planetary-atmospheres:7}
(a)~Em quatro bilhões de anos ($\num{1.26e17}$~s), Marte perde
$\num{3.8e43}$ átomos, $\num{1.9e43}$ moléculas de água, $\qty{5.7e17}{kg}$;
sobre $4\pi(\num{3.3895e6})^2 = \qty{1.44e14}{m^2}$, o que corresponde a uma camada
de \qty{3.9}{m} de espessura. (b)~Vinte vezes mais, \qty{78}{m}. (c)~Mesmo a
taxa de alta obliquidade atinge apenas o extremo inferior do intervalo de 100 a
\qty{1500}{m}; o restante foi para a crosta e o gelo, ou foi perdido
mais rapidamente quando o Sol era jovem.
\end{solution}

\begin{solution}{exo:b3:planetary-atmospheres:8}
(a)~$m = pA/g = 636\times\num{1.44e14}/3.71 = \qty{2.5e16}{kg}$.
(b)~$0.1\times\num{1.26e17} = \qty{1.3e16}{kg}$, metade da
atmosfera atual; toda ela em $\num{2.5e16}/0.1 = \qty{2.5e17}{s}$,
\num{7.8}~bilhões de anos. (c)~O Sol jovem era mais ativo e suas
tempestades mais frequentes, e a perda cresce durante as tempestades; outros canais
(hidrogênio, pulverização por íons, escape fotoquímico) não são contabilizados.
\end{solution}

\begin{solution}{exo:b3:planetary-atmospheres:9}
(a)~$S(1 - 0.193)/4 = 291$ dá $S = \qty{1442}{W.m^{-2}}$, 1.06 vezes
a da Terra. (b)~Vênus hoje absorveria $0.807\times2601/4 =
\qty{525}{W.m^{-2}}$; um Vênus jovem, $0.7\times2601 = \qty{1821}{W.m^{-2}}$
recebido, \qty{367}{W.m^{-2}} absorvido: ambos acima do limite, de modo que um
Vênus jovem sem nuvens teria entrado em \omterm{def:b3:planetary-atmospheres:runaway}{estufa descontrolada} desde o início. (c)~As nuvens
elevam o albedo; modelos tridimensionais com nuvens encontraram um caminho
para um Vênus jovem habitável.
\end{solution}

\begin{solution}{exo:b3:planetary-atmospheres:10}
(a)~$H = kT/mg = \num{1.381e-23}\times1000/(16\times\num{1.661e-27}\times9.21)
= \qty{56}{km}$. (b)~$n_\mathrm{exo} = 1/(\sigma_c H) = 1/(\num{3e-19}\times
\num{5.64e4}) = \qty{5.9e13}{m^{-3}}$; $z = 200 + 56.4\ln(\num{1e16}/
\num{5.9e13}) = 200 + 56.4\times5.13 = \qty{490}{km}$. (c)~A gravidade mais fraca
nas altitudes mais altas faz a altura de escala crescer com a altitude, de modo que a densidade cai
mais lentamente do que o assumido: a \omterm{def:b3:planetary-atmospheres:exobase}{exobase} sobe.
\end{solution}

\begin{solution}{exo:b3:planetary-atmospheres:11}
(a)~$\Phi = \num{3.0e13}\times\num{3e-3} = \qty{9e10}{cm^{-2}.s^{-1}} =
\qty{9e14}{m^{-2}.s^{-1}}$; sobre $4\pi(\num{6.371e6})^2 =
\qty{5.10e14}{m^2}$, $\num{4.6e29}$ moléculas de $\ce{H2}$ por segundo.
(b)~Um oceano contém $\num{1.4e24}/18.015\times\num{6.022e23} =
\num{4.7e46}$ moléculas de água, uma de $\ce{H2}$ para cada: ela dura
$\qty{1.0e17}{s} = \num{3.2}\text{ bilhões de anos}$, comparável à idade da
Terra, o que explica por que uma razão de mistura de \num{3e-3} marca o limite da
\omterm{def:b3:planetary-atmospheres:runaway}{estufa úmida}.
\end{solution}

\begin{solution}{exo:b3:planetary-atmospheres:12}
(a)~$p_1 = g/D\kappa = 9.81/(1.66\times0.005) = \qty{1180}{Pa}$.
(b)~$1/T = 1/273.16 - \ln(1182/611.657)\times461.5/\num{2.5e6}$, $T =
\qty{282.5}{K}$, e $\sigma T^4 = \qty{361}{W.m^{-2}}$. (c)~O limite real
é menor, portanto o nível radiante é mais frio e está a uma pressão menor: o vapor de água real
é mais opaco que este gás cinza, e seu $\kappa$ efetivo
é maior.
\end{solution}

\begin{solution}{pb:b3:planetary-atmospheres:1}
\textbf{1.} $g = GM/R^2 = \num{6.674e-11}\times\num{5.2e24}/(\num{6.2e6})^2 =
\qty{9.03}{m.s^{-2}}$; $\bar\rho = M/\tfrac43\pi R^3 = \qty{5210}{kg.m^{-3}}$.

\textbf{2.} $v_\mathrm{esc} = \sqrt{2GM/r} = \sqrt{2\times\num{3.47e14}/\num{6.7e6}}
= \qty{10.18}{km.s^{-1}}$.

\textbf{3.} $u = \sqrt{2kT/m}$: hidrogênio \qty{4.81}{km.s^{-1}}, hélio
\qty{2.41}{km.s^{-1}}, oxigênio \qty{1.21}{km.s^{-1}}.

\textbf{4.} $\lambda = (v_\mathrm{esc}/u)^2$: 4.49, 17.8, 71.2.

\textbf{5.} Só o hidrogênio: sua fração de escape é 0.06, a do hélio
\num{3.5e-7}, a do oxigênio \num{9e-30}. Nenhuma está perto do regime hidrodinâmico,
$\lambda \approx 1.5$ a 3.

\textbf{6.} $\lambda = 4.49\times1400/3000 = 2.09$: dentro da transição;
o hidrogênio teria escapado hidrodinamicamente, talvez arrastando
gases mais pesados.

\textbf{7.} $(1 + 4.49)\,e^{-4.49} = 0.0618$.

\textbf{8.} $n u/2\sqrt\pi = \num{2e11}\times4806/3.545 = \qty{2.71e14}{m^{-2}.s^{-1}}$;
$\Phi_\mathrm J = \num{2.71e14}\times0.0618 = \qty{1.68e13}{m^{-2}.s^{-1}}$.

\textbf{9.} $4\pi(\num{6.7e6})^2 = \qty{5.64e14}{m^2}$; $\num{9.5e27}$ átomos por
segundo.

\textbf{10.} $\num{1.4e24}/18.015 = \qty{7.8e22}{mol}$, $\num{4.68e46}$ moléculas,
$\num{9.36e46}$ átomos de hidrogênio.

\textbf{11.} $\num{9.36e46}/\num{9.5e27} = \qty{9.9e18}{s}$, cerca de 310 bilhões
de anos.

\textbf{12.} $\lambda_\mathrm D = 4.49\times2.014/1.008 = 8.96$; fração de escape
$9.96\,e^{-8.96} = \num{1.28e-3}$; eficiência relativa $0.707\times
\num{1.28e-3}/0.0618 = 0.015$.

\textbf{13.} $H_a = kT/m_a g = \num{1.381e-23}\times200/(44.01\times\num{1.661e-27}
\times9.03) = \qty{4.18}{km}$.

\textbf{14.} $b = \qty{1.2e21}{m^{-1}.s^{-1}}$, logo
\[
\Phi_\mathrm{DL} = \frac{b f(m_a - m_i)}{m_a H_a}
= \frac{\num{1.2e21}\times\num{1e-3}\times(42.0/44.0)}{4185}
= \qty{2.74e14}{m^{-2}.s^{-1}}.
\]

\textbf{15.} $\num{3.0e13}\times\num{1e-3}\times9.03/9.81 =
\qty{2.76e10}{cm^{-2}.s^{-1}} = \qty{2.76e14}{m^{-2}.s^{-1}}$: o mesmo
dentro de 1\,\%.

\textbf{16.} $4\pi(\num{6.2e6})^2 = \qty{4.83e14}{m^2}$; $\num{1.32e29}$ moléculas
de $\ce{H2}$, $\num{2.64e29}$ átomos por segundo.

\textbf{17.} O limite de difusão fornece 28 vezes o que a \omterm{def:b3:planetary-atmospheres:exobase}{exobase} deixa escapar com
sua densidade atual. O hidrogênio se acumula na termosfera até que o
fluxo de Jeans iguale o suprimento: a longo prazo, o escape é
limitado por difusão, como na Terra.

\textbf{18.} $\num{2e11}\times28 = \qty{5.6e12}{m^{-3}}$.

\textbf{19.} $\num{4.68e46}/\num{1.32e29} = \qty{3.54e17}{s}$, 11.2 bilhões de anos.

\textbf{20.} $S = 0.30/0.60^2 = 0.833$.

\textbf{21.} $d = \sqrt{L/S_\mathrm{eff}}$: borda interna $\sqrt{0.30/0.939} =
\qty{0.565}{au}$, borda externa $\sqrt{0.30/0.290} = \qty{1.02}{au}$. Kepos,
a \qty{0.60}{au}, está dentro, perto da borda interna.

\textbf{22.} $0.939/0.833 = 1.13$: um aumento de brilho de 13\,\%.

\textbf{23.} $0.833\times1361\times0.7/4 = \qty{198}{W.m^{-2}}$, $T_\mathrm{eq} =
(198/\sigma)^{1/4} = \qty{243}{K}$; bem abaixo do
\qty{291}{W.m^{-2}} limite (calculado para uma estrela semelhante ao Sol, mas a margem
é ampla).

\textbf{24.} Três vezes mais rápido: 3.7 bilhões de anos.

\textbf{25.} \textbf{11~bilhões~de~anos}.
\end{solution}
